\section{Sequence models and signal hypothesis testing}\label{sec:seq}

Section~\ref{sec:ml} used a small convolutional network. Two natural questions follow, both asked on the same data, the same split by chunk, the same test injections and the same thresholds on the validation noise: (i) does a structure designed for sequences --- a recurrent network or a Transformer~\cite{Vaswani} --- detect bursts better than the CNN, and why or why not; (ii) for a signal that is above the detection threshold and incompatible with the noise, which method can tell \emph{which morphology} produced it?
Because the AUC of the CNN ensemble is already $\approx0.99$, the detectors are judged by the efficiency at a fixed low false-alarm probability ($\snr_{50}$, $\snr_{90}$ at FAP $10^{-2}$ and $10^{-3}$) and by their behaviour on glitches, not by the AUC.

\subsection{Transformer and GRU against the CNN}
\paragraph{Architectures.} The 1024 samples are cut into 64 non-overlapping patches of 16 samples (15.6\,ms)~\cite{ViT} and embedded linearly ($d=64$). The \emph{Transformer} has a learned positional embedding, three pre-norm blocks (4 heads, MLP 128, dropout $0.1$) and a head that concatenates max- and average-pooling over the tokens ($1.1\times10^5$ parameters). The \emph{GRU}~\cite{GRU} has the same patch embedding and head, and two stacked bidirectional GRU layers of 48 units ($8.8\times10^4$ parameters): it isolates the sequence mixer (recurrence against attention). The CNN of Section~\ref{sec:ml} has $1.8\times10^4$ parameters.
All are trained with AdamW~\cite{AdamW}, two epochs of warm-up and cosine decay, batch 128, early stopping on the validation loss, at most 12 epochs (CPU budget); the first seed of the Transformer was also trained for 25 epochs and ends with the same validation loss ($0.312$ against $0.310$).
\paragraph{Pre-training.} A masked-patch self-supervised pre-training~\cite{MAE} (40\,\% of the patch embeddings replaced by a mask token, a linear decoder reconstructs the raw samples) on the \emph{noise-only} training windows is tested as a way to save labelled injections.

\paragraph{Results.} Table~\ref{tab:seq} and Fig.~\ref{fig:seqeff} give the efficiency. The CNN ensemble and the Bayes factor of the three families are the best at FAP $10^{-2}$; the GRU and the Transformer are behind the CNN by $0.7$ to $2.1$ SNR units (the smallest gap is on pulses, the largest on chirps for the Transformer); at FAP $10^{-3}$ the ranking is unchanged and the gap grows (on chirps the Transformer needs $\snr_{50}=8.8$ against $6.2$). Four further observations:
\begin{itemize}[nosep]
\item \textbf{The number of parameters is not the explanation.} The Transformer has six times the parameters of the CNN. The CNN's local filters followed by a global max-pool encode the two facts the task needs --- a burst is localised and can be anywhere in the window --- and at this training-set size ($\sim1.6\times10^4$ windows per class) that prior is worth more than the flexibility of attention. The label-efficiency curve (Fig.~\ref{fig:label}) agrees: with 1000 labelled windows per class the CNN is at AUC $0.91$ and the Transformer at $0.54$.
\item \textbf{Pre-training learns almost nothing.} The masked-patch error is $0.995$ of the error of the trivial predictor ``zero'': after whitening the noise is nearly white, so there is little to predict from the neighbours (the same reason why the autoencoder of Section~\ref{sec:ml} does not beat excess power). The pre-trained encoder is a better starting point only with 1000 labels per class (AUC $0.76$ against $0.54$) and equal to the random start with 4000 labels or the full set; at the level of $\snr_{50}$ it changes nothing ($6.29/4.87/7.29$ against $6.38/4.88/7.35$).
\item \textbf{The attention is informative} (Fig.~\ref{fig:att}). The most attended patch is within $\pm1$ patch of the burst time in $48\,\%$ (sine-Gaussian), $61\,\%$ (pulse) and $44\,\%$ (chirp) of the windows, against $4.7\,\%$ by chance; chirps, which last 0.1--0.2\,s, are the least localised and are where the Transformer is weakest.
\item \textbf{Glitch rejection is not better.} Taking the $5\,\%$ of test noise windows with the largest $\max|v|$, the alarm rate inside this subset at the threshold of FAP $10^{-2}$ is $2.7\,\%$ (CNN ensemble), $4.9\,\%$ (Transformer), $5.2\,\%$ (GRU), $5.7\,\%$ (Bayes) and $6.9\,\%$ (bank); the efficiency for $\snr=8$ injections is $0.96$--$0.98$ (CNN), $0.83$ (Transformer), $0.995$ (Bayes). A glitch does not hide a $\snr=8$ signal for any method (the efficiency inside the glitchy windows equals the overall one).
\end{itemize}

\begin{table}[h]\centering\small
\begin{tabular}{lrrrrrr}
\toprule
& \multicolumn{3}{c}{FAP $10^{-2}$: $\snr_{50}$} & \multicolumn{3}{c}{FAP $10^{-3}$: $\snr_{50}$}\\
statistic & SG & pulse & chirp & SG & pulse & chirp\\
\midrule
matched-filter bank (SG templates) & 4.87 & 5.34 & 6.28 & 6.80 & 7.37 & 9.01\\
Bayes factor (3 families) & 5.21 & 4.78 & 5.34 & 7.56 & 7.24 & 7.72\\
CNN ensemble (5) & \textbf{5.05} & \textbf{4.21} & \textbf{5.28} & \textbf{6.09} & \textbf{5.04} & \textbf{6.24}\\
Transformer (2 seeds) & 6.38 & 4.88 & 7.35 & 7.46 & 5.73 & 8.78\\
Transformer, pre-trained encoder & 6.29 & 4.87 & 7.29 & 7.26 & 5.73 & 8.77\\
GRU & 6.27 & 5.16 & 6.36 & 7.44 & 6.05 & 7.70\\
\bottomrule
\end{tabular}
\caption{Optimal SNR at 50\,\% efficiency of the sequence models, the CNN and the model-based statistics (thresholds on the validation noise; test injections). Bold: best neural network in each column.}\label{tab:seq}
\end{table}

\begin{figure}[h]\centering
\includegraphics[width=0.95\textwidth]{fig06_efficiency.pdf}
\caption{Efficiency against optimal SNR at FAP $10^{-2}$ (top) and $10^{-3}$ (bottom) for the model-based statistics, the CNN ensemble, the Transformer (with and without the pre-trained encoder) and the GRU.}\label{fig:seqeff}
\end{figure}

\begin{figure}[h]\centering
\includegraphics[width=0.55\textwidth]{fig06_label_efficiency.pdf}
\caption{Test AUC against the number of labelled windows per class for the CNN, the Transformer and the Transformer with the pre-trained encoder.}\label{fig:label}
\end{figure}

\begin{figure}[h]\centering
\includegraphics[width=0.95\textwidth]{fig06_attention.pdf}
\caption{Attention received by each patch (green) in the last block of the Transformer for one $\snr=10$ injection of each morphology; dashed: burst time.}\label{fig:att}
\end{figure}

\subsection{Signal hypothesis testing: which morphology?}
Given a signal above threshold, the decision among the three families is compared across methods (Fig.~\ref{fig:morph}). The frequentist rules use the maximum matched-filter statistic of each family (a bank of 160 sine-Gaussian, 14 pulse and 120 chirp templates), either raw or converted into a background p-value to take the different trial factors into account; the Bayesian rule uses the marginal likelihood (Bayes factor) of each family with $\sigma_\rho=10$ and a uniform prior over the templates, so that the Occam factor penalises the larger families; the networks are trained as four-class classifiers (noise, sine-Gaussian, pulse, chirp) on the same data as Section~\ref{sec:ml}.
All methods use the same rule for the four-class decision: a window is called noise if its own detection statistic is below the FAP $10^{-2}$ threshold of the validation noise, otherwise the family with the largest score is chosen.

\begin{table}[h]\centering\small
\begin{tabular}{lrrrrr}
\toprule
3-way accuracy at $\snr=$ & 5 & 6 & 8 & 10 & 14\\
\midrule
Bayes (3 families) & 0.771 & \textbf{0.906} & \textbf{0.969} & 0.977 & 0.989\\
frequentist, max $\rho^2$ & 0.724 & 0.837 & 0.921 & 0.948 & 0.981\\
frequentist, background p-value & \textbf{0.825} & 0.900 & 0.913 & 0.889 & 0.816\\
CNN 4-class (2 seeds) & 0.725 & 0.863 & 0.960 & \textbf{0.980} & \textbf{0.991}\\
GRU 4-class & 0.665 & 0.813 & 0.925 & 0.969 & 0.989\\
Transformer 4-class & 0.643 & 0.759 & 0.867 & 0.925 & 0.955\\
Transformer 4-class, pre-trained & 0.665 & 0.765 & 0.873 & 0.924 & 0.955\\
\bottomrule
\end{tabular}
\caption{Accuracy of the choice among the three families for injections of each SNR (chance: 0.333). The p-value row is plotted in Fig.~\ref{fig:morph}; bold: best in each column (at $\snr=10$ and 14 the CNN and the Bayes factor differ by less than the statistical error, $\approx0.01$).}\label{tab:morph}
\end{table}

\begin{itemize}[nosep]
\item \textbf{Bayesian model comparison is the best classifier from $\snr\approx6$ among the networks trained on the static set (tied with the CNN above $\snr\approx10$)} ($0.969$ at $\snr=8$; this ranking changes when the networks are trained with fresh injections, Section~\ref{sec:tuning}) and has the best four-class balanced accuracy in both SNR ranges ($0.778$ for $\snr=4$--7, $0.988$ for $\snr=8$--16). The CNN is close at high SNR ($0.960$ at $\snr=8$) and $0.04$--$0.05$ behind at $\snr=5$--6; then come the GRU, and last the Transformer ($0.867$--$0.873$), pre-trained or not. The answer to the question is therefore: among the networks the convolutional structure (local filters plus global pooling) does best, a bidirectional recurrent encoder is second and attention on 16-sample patches third; none beats the Bayesian comparison, which uses the exact likelihood of each family.
\item \textbf{The frequentist rule depends on how the trial factors are treated.} The largest raw $\rho^2_f$ is $0.05$ worse than the Bayes factor at $\snr=8$; converting each into a background p-value is the best rule at $\snr\le5$ but its accuracy \emph{falls} with SNR beyond $\snr\approx9$ ($0.82$ at 14) because the empirical p-values of loud signals saturate at $1/(N+1)$ for every family and the tie is broken by the raw $\rho^2$, which favours the families with more templates. This is a limit of the $8112$-window background, not of the method.
\item \textbf{Detection comes first.} Only $\approx46\,\%$ of the $\snr=5$ injections and $97\,\%$ of those at $\snr=7$ pass the common Bayesian gate, so the morphology can be discussed only for signals loud enough to be detected. In the range $\snr=4$--7 the fraction of signals called noise is $0.40$ (Bayes, frequentist), $0.51$ (CNN), $0.54$ (GRU), $0.60$ (Transformer).
\item \textbf{A signal incompatible with every hypothesis.} A Gaussian-enveloped burst of band-limited stochastic noise ($\sigma_t=15$\,ms) belongs to none of the families. It is detected when loud ($94$--$98\,\%$ at $\snr=12$) and every method gives it a family, mostly sine-Gaussian ($0.93$ max-$\rho^2$, $0.95$ CNN, $0.80$ Bayes, $0.82$--$0.85$ Transformers, $0.90$ GRU) and then pulse. No method can say ``none of the above'': a softmax is normalised over the classes it knows and the Bayesian posterior over the families it contains. Flagging such signals would need an explicit goodness-of-fit (the residual after subtracting the best template) or an open-set head; neither was tried.
\end{itemize}

\begin{figure}[h]\centering
\includegraphics[width=0.95\textwidth]{fig06_morphology_accuracy.pdf}
\caption{Accuracy of the morphology choice against SNR: all injections (left), injections above the common Bayesian detection threshold (centre) and the fraction passing it (right).}\label{fig:morph}
\end{figure}

\subsection{Is the Transformer's loss to the CNN real? A tuning study}\label{sec:tuning}
The comparison above used a small budget (at most 12 epochs, two seeds, no tuning of depth, width or patch size), so ``the Transformer loses'' is a statement about that budget. It is stress-tested by changing \emph{one lever at a time}, with the same data, thresholds and test injections (verified identical to those of the previous sections), and by giving the CNN the same levers so that the comparison stays fair.

\paragraph{Diagnosis.} At the best-validation epoch every model has a small train--validation gap ($0.01$--$0.03$), but the \emph{level} of the loss differs: $0.15$--$0.17$ for the CNN of Section~\ref{sec:ml}, $0.29$--$0.31$ for the Transformers and $0.29$ for the GRU (Fig.~\ref{fig:tuning}, left). The sequence models are limited by what they can fit, not by over-fitting. The CNN trained with the same loop on the static data does the opposite: it over-fits (training loss $0.07$--$0.09$, validation $0.17$--$0.19$, and a validation loss that rises again after epoch 15). So the CNN is \emph{data-limited}, the patch Transformer is \emph{capacity/bias-limited}.

\paragraph{Levers.} (i) \emph{Data}: fresh injections --- new morphology, SNR, parameters and phase --- at every epoch, for the CNN, the Transformer and the hybrid; (ii) \emph{optimisation}: peak learning rate $3\times10^{-3}$ instead of $10^{-3}$, and the 25-epoch run; (iii) \emph{resolution}: patches of 8 samples (128 tokens); (iv) \emph{inductive bias}: a \emph{convolutional stem} (the backbone of the CNN, down-sampling by 8) followed by two Transformer blocks on 128 tokens (the ``hybrid''). Budget: 12 epochs for the sequence models, 20 for the CNN; the Transformer variants take 23--81\,min each on a laptop.

\begin{table}[h]\centering\small
\begin{tabular}{lrrrr}
\toprule
model & train loss & val.\ loss & $\overline{\snr_{50}}$ @$10^{-2}$ & $\overline{\snr_{50}}$ @$10^{-3}$\\
\midrule
Transformer (Sec.~\ref{sec:seq}, 2 seeds) & 0.30 & 0.31 & 6.20 & 7.32\\
\quad + fresh injections (2 seeds) & 0.30 & 0.30 & 6.27 & 7.38\\
\quad + learning rate $3\times10^{-3}$ & 0.30 & 0.32 & 6.36 & 7.68\\
\quad + patches of 8 samples & 0.31 & 0.31 & 6.59 & 7.58\\
\midrule
hybrid (conv.\ stem + 2 blocks), static (2 seeds) & 0.22 & 0.27 & 5.64 & 6.79\\
hybrid, fresh injections (2 seeds) & 0.23 & 0.25 & 5.33 & 6.21\\
\midrule
CNN, static data (2 seeds) & 0.08 & 0.18 & 5.22 & 6.26\\
CNN, fresh injections (2 seeds) & 0.09 & \textbf{0.117} & \textbf{4.60} & \textbf{5.65}\\
CNN ensemble of Sec.~\ref{sec:ml} (5, static) & -- & -- & 4.85 & 5.79\\
\bottomrule
\end{tabular}
\caption{Tuning study: training and validation loss at the best epoch and mean over the three morphologies of $\snr_{50}$ at FAP $10^{-2}$ and $10^{-3}$ (thresholds on the validation noise). Differences below $\approx0.3$ are not significant: the half-range between seeds of the same configuration is $0.11$/$0.07$ (static CNN), $0.03$/$0.01$ (fresh CNN), $0.03$/$0.13$ (Transformer, fresh), $0.02$/$0.18$ (hybrid, static) and $0.03$/$0.07$ (hybrid, fresh) at FAP $10^{-2}$/$10^{-3}$; the lr-$3\times10^{-3}$ and patch-8 variants have a single seed.}\label{tab:tuning}
\end{table}

\begin{figure}[h]\centering
\includegraphics[width=0.95\textwidth]{fig07_training_curves.pdf}\\[0.3em]
\includegraphics[width=0.95\textwidth]{fig07_levers.pdf}
\caption{Top: training and validation loss of the variants. Bottom: mean $\snr_{50}$ over the three morphologies for every model, at FAP $10^{-2}$ and $10^{-3}$ (dotted: the Transformer of Section~\ref{sec:seq}).}\label{fig:tuning}\label{fig:levers}
\end{figure}

\paragraph{Findings.}
\begin{enumerate}[nosep]
\item \textbf{The patch Transformer does not recover with any single lever.} Fresh injections ($6.27$, two seeds), a three times higher learning rate, smaller patches (four times the cost) and the 25-epoch run leave $\overline{\snr_{50}}$ at $6.3$--$6.6$ and the validation loss at $0.30$--$0.32$: the limitation is the embedding, not the data or the optimisation time. Chirps are the worst case ($\snr_{50}\approx7.4$--$8.0$ against $5.5$--$5.8$ for the CNN): a 30--80\,ms oscillation needs frequency-selective filters, which a linear map of 16 raw samples does not provide.
\item \textbf{A convolutional stem is the lever that works.} The hybrid reaches $5.64$ (static) and $5.33$ (fresh), $0.6$--$0.9$ better than the patch Transformer, well above the seed scatter, and a validation loss of $0.27$/$0.25$; it is the only sequence model competitive at FAP $10^{-3}$ ($6.21$ against $7.3$--$7.7$).
\item \textbf{But the CNN stays ahead with the same data.} Static data: hybrid $5.64$ against CNN $5.22$; fresh data: $5.33$ against $4.60$. Attention on top of a good convolutional front-end does not add detection power on one-second windows with a localised burst.
\item \textbf{The CNN of Section~\ref{sec:ml} was data-limited, and fresh injections are the best lever of the whole study.} Fresh injections cut $\overline{\snr_{50}}$ of the CNN from $5.22$ to $4.60$ at FAP $10^{-2}$ and from $6.26$ to $5.65$ at $10^{-3}$, of the hybrid by $0.31$ and $0.58$, and of the patch Transformer by nothing. A \emph{single} fresh CNN is better than the five-member static ensemble ($4.85$, $5.79$), with the same $\snr_{90}$ ($6.6$); its validation loss falls from $0.17$--$0.19$ to $0.117$. This is the reference CNN for any future comparison; the numbers of Section~\ref{sec:ml} are unchanged but should be read as those of a data-limited network.
\item \textbf{The comparison does not favour the CNN.} At a nominal FAP of $1\,\%$ the observed false-alarm probability on the test chunks is $0.4$--$0.5\,\%$ for the fresh CNNs, $0.7$--$1.0\,\%$ for the patch Transformers and $1.0$--$1.1\,\%$ for the hybrids: the CNN threshold is the more conservative one, so at equal observed FAP the gap would be larger.
\end{enumerate}

\paragraph{Intuition (interpretation, not tested separately).} The burst is a short, localised oscillation in white noise; seeing it requires a filter tuned to its frequency, which is what a convolution does (a learned matched filter). A patch Transformer first mixes 16 raw samples with a single linear map, which cannot build such filters, and nothing downstream recovers what is lost: this is consistent with the high training loss. Once the filters are provided by a convolutional stem the Transformer improves a lot, but a burst that sits in one place of a one-second window leaves little long-range structure for attention to exploit, and recurrence, whose strength is remembering distant context, is similarly wasted.

\paragraph{Multiclass networks with fresh injections change a conclusion of Section~\ref{sec:seq}.} The four-class networks of the previous subsection were trained on the static set. Trained with fresh injections (Table~\ref{tab:mcfresh}) the CNN rises from $0.960$ to $0.977$ three-way accuracy at $\snr=8$ and from $0.983$ to $0.991$ balanced four-class accuracy for $\snr=8$--16 ($0.796$ against $0.722$ for $\snr=4$--7): it ties with the Bayesian model comparison at low SNR ($0.768$ against $0.771$ at $\snr=5$) and is slightly \emph{ahead} of it at high SNR ($0.977$ against $0.969$ at $\snr=8$, $0.988$ against $0.977$ at $\snr=10$; differences of $\approx0.01$, about two standard errors with 1500 injections per point). The hybrid network is level with the Bayes factors; the patch Transformer is not helped ($0.844$ at $\snr=8$ against $0.867$ static). The statement of Section~\ref{sec:seq} that Bayesian model comparison is the best morphology classifier was therefore a statement about a data-limited CNN: with the same budget and fresh injections the CNN is as good.

\begin{table}[h]\centering\small
\begin{tabular}{lrrrrrcc}
\toprule
3-way accuracy at $\snr=$ & 5 & 6 & 8 & 10 & 14 & bal.\ 4--7 & bal.\ 8--16\\
\midrule
Bayes (3 families) & 0.771 & 0.906 & 0.969 & 0.977 & 0.989 & 0.778 & 0.988\\
CNN 4-class, static & 0.725 & 0.863 & 0.960 & 0.980 & 0.991 & 0.722 & 0.983\\
CNN 4-class, fresh & 0.768 & 0.900 & \textbf{0.977} & \textbf{0.988} & \textbf{0.997} & \textbf{0.796} & \textbf{0.991}\\
hybrid 4-class, fresh & 0.769 & 0.888 & 0.965 & 0.986 & 0.997 & 0.760 & 0.987\\
Transformer 4-class, static & 0.643 & 0.759 & 0.867 & 0.925 & 0.955 & 0.661 & 0.944\\
Transformer 4-class, fresh & 0.635 & 0.736 & 0.844 & 0.899 & 0.950 & 0.656 & 0.933\\
GRU 4-class, static & 0.665 & 0.813 & 0.925 & 0.969 & 0.989 & 0.700 & 0.966\\
\bottomrule
\end{tabular}
\caption{Three-way morphology accuracy at several SNRs and balanced four-class accuracy for $\snr=4$--7 and 8--16, with static and fresh-injection training (the Bayes row is from Table~\ref{tab:morph}).}\label{tab:mcfresh}
\end{table}

\paragraph{Caveats.} A single seed for the lr-$3\times10^{-3}$ and patch-8 Transformers and for the four-class Transformer and hybrid; a budget of 12 epochs (the hybrid reaches its best validation epoch at 7--12, so the fresh-data runs might gain a little more); depth, width, number of heads and the pre-training were not swept; the unit tested is one second of whitened data with one localised burst --- longer windows, or several overlapping signals, are the setting where attention could be expected to matter.

\subsection{Seed scatter and paired tests}\label{sec:seeds}
The comparisons above used one or two seeds per architecture. Notebook~30 trains more seeds with the same recipe (CNN with fresh injections: $5$; hybrid: $3$; Transformer: $3$) and evaluates all on the same test injections with the validation thresholds (mean $\snr_{50}$ over the three morphologies).
\begin{table}[h]\centering\small
\begin{tabular}{lrrr}
\toprule
architecture (seeds) & FAP $10^{-2}$: mean $\pm$ std (range) & FAP $10^{-3}$: mean $\pm$ std\\
\midrule
CNN, fresh injections (5) & $4.31\pm0.40$ ($3.72$--$4.63$) & $5.67\pm0.03$\\
hybrid, fresh (3) & $5.35\pm0.05$ & $6.29\pm0.16$\\
Transformer, fresh (3) & $6.29\pm0.05$ & $7.47\pm0.20$\\
average of the 5 CNN logits & $3.79$ & $5.32$\\
\bottomrule
\end{tabular}
\caption{Seed-to-seed scatter of the detectors (the test-set bootstrap error of a single network is $0.08$).}\label{tab:seeds}
\end{table}
The scatter of the CNN at FAP $10^{-2}$ ($0.40$; two seeds at $3.72$ and $4.08$ against $4.56$--$4.63$ for the other three) is much larger than the $0.03$ half-range of the two seeds of notebook~07 and larger than the test-set error, while at $10^{-3}$ all seeds agree to $0.03$: a single-seed number at $10^{-2}$ is uncertain by about $\pm0.4$. The architecture differences are nevertheless larger than the scatter and significant at the level of seeds (permutation and Welch tests): hybrid minus CNN $+1.04$ ($p=0.018$ permutation, $0.004$ Welch) and Transformer minus CNN $+1.98$ ($p=0.018$, $0.0003$) at $10^{-2}$; $+0.62$ and $+1.79$ at $10^{-3}$. Averaging five CNNs lowers $\snr_{50}$ by a further $0.5$ ($3.79$, $5.32$), the largest single gain obtained in this study at fixed training data. The earlier statement that the ranking CNN $<$ hybrid $<$ Transformer is robust is confirmed; its margin is $1$ to $2$ units, not tenths.
\begin{figure}[h]\centering\includegraphics[width=0.75\textwidth]{fig30_seed_scatter.pdf}\caption{Mean $\snr_{50}$ of every seed (dots) with mean and standard deviation (bar).}\label{fig:seeds}\end{figure}

\subsection{Other architectures: ResNet-1D, TCN, spectrogram CNN, graph network}\label{sec:otherarch}
Notebook~34 trains four more detectors with the recipe of notebook~07 (fresh injections, $12$ epochs, one seed): a ResNet-1D (stem and four residual stages), a temporal convolutional network (residual dilated convolutions, dilation $1$--$32$), a 2-D CNN on the log-spectrogram ($33\times121$) and a graph network (EdgeConv with $8$ nearest neighbours) on the $64$ time-frequency pixels of largest power, the construction behind the clustering of coherent WaveBurst. Mean $\snr_{50}$ over the three morphologies at FAP $10^{-2}$ ($10^{-3}$): ResNet-1D $5.01$ ($5.77$), TCN $5.09$ ($5.87$), 2-D CNN $5.20$ ($5.97$), graph network $5.48$ ($6.46$), against $4.31\pm0.40$ ($5.67\pm0.03$) for the project CNN over five seeds, $5.35$ ($6.29$) for the hybrid and $6.29$ ($7.47$) for the Transformer; the bootstrap error of each is $0.04$--$0.05$ at $10^{-2}$ and the seed scatter of the CNN is $0.4$. No architecture beats the plain 1-D CNN, the four are within $\approx0.2$--$0.5$ of each other and of the matched-filter bank ($5.5$), and all are better than the patch Transformer: the global max-pooling over localised convolutional features is what the detection of a localised burst needs, and adding depth, dilation or a second dimension does not change that. The graph network, which sees only $64$ pixels, loses $\approx0.5$ against the best and is the cheapest to evaluate; per morphology the four differ by less than $1$ ($4.5$--$5.8$ over the three morphologies), without a consistent winner.
\begin{figure}[h]\centering\includegraphics[width=0.75\textwidth]{fig34_architectures.pdf}\caption{Mean $\snr_{50}$ of the four additional architectures (blue) and of the references of Section~\ref{sec:seeds} (green); light part: increase at FAP $10^{-3}$.}\label{fig:otherarch}\end{figure}

\subsection{Transformer tuning, second round: depth, width, length}\label{sec:tftune2}
The levers left open by notebook~07 were run in notebook~39 (patch $16$, fresh injections, one seed each): six blocks instead of three ($214\,$k parameters), width $128$ with MLP $256$ ($425\,$k) and $30$ epochs instead of $12$. Mean $\snr_{50}$ at FAP $10^{-2}$ ($10^{-3}$): baseline $6.29$ ($7.47$; three seeds, seed scatter $0.05$ ($0.2$)), depth $6.23$ ($7.21$), width $6.11$ ($6.97$), $30$ epochs $5.81$ ($7.15$), with validation losses $0.301$, $0.299$, $0.301$ and $0.255$. Deeper and wider networks gain $0.06$--$0.2$ at $10^{-2}$ (within the seed scatter for the depth) and $0.3$--$0.5$ at $10^{-3}$; training longer is the only lever whose gain at $10^{-2}$ ($0.48$) is far beyond the scatter, and it still leaves the Transformer $1.5$ units behind the CNN ($4.31$) and $0.5$ behind the hybrid ($5.35$). Combined with the earlier levers (learning rate, patch size, fresh data, convolutional stem), the Transformer on one-second windows with a localised burst does not reach the CNN with any budget tried here; contrastive pre-training and overlapping windows, also on the original list, were not run.
\begin{figure}[h]\centering\includegraphics[width=0.98\textwidth]{fig39_transformer_tuning.pdf}\caption{Left: validation loss of the variants. Right: mean $\snr_{50}$ of the variants and of the references.}\label{fig:tftune2}\end{figure}
